% =====================================================================
% §3 SETTING AND DATA — restructured
% Corpus pattern: institutional background -> data -> stylized facts,
% all before the model. Assembled programmatically: 3.1 and 3.2's framing
% prose is new; the institutional paragraph, the data subsections, and the
% four stylized facts are relocated verbatim from elsewhere in the paper.
%
% WIRING DEPENDENCIES (all in the same commit that inputs this file):
%   1. delete the old \section{Setting and Stylized Facts} block
%   2. delete the old \section{Data} block  (its \label{sec:data} is
%      carried onto §3.2 below -- 4 live \ref{sec:data} calls depend on it)
%   3. delete the institutional paragraph from the identification section
%      ("Before July 2021, the French DPE was a paper exercise...")
% =====================================================================

\section{The French Energy-Performance Disclosure Regime}
\label{sec:setting}

This section sets out the institution, the data, and four features of the market that any
candidate theory of EPC disclosure has to accommodate. We lead with apartments throughout and
treat houses as a corroborating segment. The two display the same qualitative patterns, but the
apartment market has the lower audit coverage, the larger share of strategic concealment, and the
stability across the reform that the identification strategy exploits.

\subsection{Institutional Background}
\label{ss:institutions}

The French \emph{diagnostic de performance \'energ\'etique} (DPE) is a standardised energy audit of
a dwelling, carried out by an accredited third-party assessor and recorded in a national registry
maintained by the \emph{Agence de la Transition \'Ecologique} (ADEME). It assigns a letter grade
from A to G summarising modelled energy consumption per square metre. Two features make it
unusual as a disclosure instrument. It is certified rather than self-reported, so a displayed
grade can be checked against the registry. And the registry is centralised and public, so the
\emph{existence} of an audit is verifiable independently of both the seller and the listing
platform. That second property is what makes the empirical exercise in this paper possible.

Before July 2021, the French DPE was a paper exercise. Scores were methodologically inconsistent across assessors, carried no legal weight, and were routinely ignored in negotiations: a buyer who later discovered a disclosed rating was misleading had no legal recourse. The reform of July 1, 2021 changed the regime fundamentally. It standardized the methodology, imposed stricter accreditation requirements on assessors, and, most importantly, rendered disclosed EPC ratings legally opposable: a buyer can now use a disclosed rating as grounds for price renegotiation or legal action. The DPE became, for the first time, a document with contractual force.

Two things the reform did not change matter as much as what it did. First, the obligation to
commission a DPE before marketing a property predates it. The requirement, and the liability
attached to marketing a property without the required diagnostic, were established by \emph{loi
ALUR} (Article L271-4 of the \emph{Code de la construction et de l'habitation}) well before 2021.
The legal exposure a seller incurs by coming to market without a certificate is therefore governed
by a mandate that straddles the reform date unchanged, while what July 2021 altered is how much a
\emph{disclosed} grade is worth to the buyer. Second, the reform sharpened a pre-existing
regulatory threshold rather than inventing one. Under \emph{loi Climat et R\'esilience}, dwellings
rated E, F, or G face a phased prohibition on new lettings. This draws a bright line between the
bottom of the grade distribution and the rest, and gives the binary quality distinction used below
a concrete regulatory anchor rather than an arbitrary cutoff.

\subsection{Data and Measurement}
\label{sec:data}

Testing a theory of strategic concealment requires observing something buyers cannot, namely whether
a silent seller holds a certificate. We recover it by matching listings against the ADEME registry.

\subsubsection{Sources and Sample Construction}

The empirical analysis draws on two sources matched at the listing level, together with transaction records used separately. Listing records are sourced from SeLoger, a major French residential real-estate platform covering sales across metropolitan France. For each listing, we observe the listing price, property characteristics ($X_i$: area, number of rooms, floor, construction era, location), the disclosure indicator $d\in\{0,1\}$, and, when $d=1$, the displayed EPC letter (A through G). Transaction records are not matched to individual listings; we use them to estimate the mapping from listing price to realized sale price ($\gamma_0$ and $\gamma_1$ in equation~\eqref{eq:payoff}), which supports treating the listing price as the relevant outcome for the disclosure decision. The ADEME registry is the official national database of certified Diagnostic de Performance Énergétique (DPE) audits, maintained by France's environment and energy management agency and reported by accredited third-party assessors. We match listings to the ADEME registry using address, property type, area, and floor, assigning to each listing an indicator $K=1$ if a certified audit record exists at the time of listing, and $K=0$ otherwise. For the structural estimation, listings for which no address is available are dropped, as the ADEME match requires a geocoded address; these account for a small share of the total sample.

The full listing sample covers 2021Q3 to 2024Q4, the post-reform period during which EPC disclosure is legally binding and DPE scores are credible to buyers. The apartment sample comprises 843,621 listings and the house sample 687,878 listings. For structural estimation we restrict to 2022 to 2024: the H2~2021 period (the six months immediately following the reform) is excluded because the market had not yet fully repriced the DPE (hedonic estimates of $\hat\beta_1$ are near zero or negative in that window), yielding 767,000 apartment listings and 633,000 house listings for the structural sample. For 51.2\% of apartment listings and 40.3\% of house listings, the ADEME match yields $K=1$; the remainder are classified as genuinely uninformed ($K=0$, no verified audit on record). For identification of market-level price indices, we use the panel of hedonic price indices constructed in \citet{patras_mandatory_2024}, covering 110 markets (attraction zones, \textit{zones d'emploi}) with monthly observations from October 2018 through July 2024.

\subsubsection{Four-Group Classification}

Combining the listing disclosure decision $d$ with the ADEME information status $K$ generates four mutually exclusive groups, summarized in Table~\ref{tab:groups}.

\begin{table}[htbp]
\centering
\caption{Classification of listings and sample composition}
\label{tab:groups}
\small\renewcommand{\arraystretch}{1.3}
\begin{tabular}{llcccc}
\toprule\toprule
\textbf{Group} & \textbf{Interpretation} & $d$ & $K$ & \textbf{Apartments} & \textbf{Houses} \\
\midrule
Truthful disclosers  & Informed seller reveals $\theta$        & 1 & 1 & 47.7\% & 39.2\% \\
Strategic hiders     & Informed seller conceals $\theta$     & 0 & 1 & \textbf{3.5\%}  & \textbf{1.1\%}  \\
Genuinely uninformed & Silent; no registry record              & 0 & 0 & 16.7\% & 8.6\% \\
Disclosed, unmatched & Displays EPC; no post-2021 registry record & 1 & 0 & 32.2\% & 51.2\% \\
\midrule
Total                &                                         &   &   & 843,621 & 687,878 \\
\bottomrule\bottomrule
\end{tabular}
\vspace{0.5em}
\begin{minipage}{\linewidth}
\footnotesize\textit{Notes.} Shares of the full listing sample of each property type, computed from the
listing--registry match. The ``disclosed, unmatched'' group displays an EPC score for which no record
appears in the public ADEME registry, which contains only certificates issued from July~2021 onward;
these listings are classified as truthful disclosers under the no-falsification assumption. Sellers
with no post-2021 registry record ($K=0$) therefore account for 48.8\% of apartment and 61.1\% of
house listings, the large majority of whom do disclose a score. Columns may not sum to exactly 100\%
because of rounding.
\end{minipage}
\end{table}

\noindent Three features stand out. First, strategic hiding is rare but economically meaningful: 3.5\% of apartment listings and 1.1\% of house listings come from informed sellers who chose not to disclose. Pooling across listings, the hiding rate among informed sellers ($K=1$) is 6.9\% for apartments and 2.7\% for houses; the corresponding unweighted means across markets, reported as $\hat\sigma_{0,m}$ in the results, are 7.7\% and 3.7\%. Second, silence is uncommon relative to non-verification: only 16.7\% of apartment listings and 8.6\% of house listings are silent sellers with no registry record, and these are the non-disclosers who provide involuntary cover for strategic hiders. Third, and quantitatively dominant, 32.2\% of apartment listings and 51.2\% of house listings display a score with no matching registry record. Taken together, sellers without a post-2021 registry record account for 48.8\% of apartment listings and 61.1\% of house listings, but the large majority of them do disclose a score.

Three features of this match bear on its credibility. It is a match to an administrative registry
rather than to a survey or a self-report, so the information-status indicator does not depend on
seller cooperation and cannot be strategically misreported. It requires a geocoded address, so the
matched sample is the subset of listings that can be located at the address level rather than the
full listing universe; sample scopes are stated alongside the statistics reported below. And, most
consequentially, $K=0$ must be read as \emph{no post-July-2021 certificate in the public registry}
rather than as \emph{no certificate exists}. The registry was populated from the reform date
onward, while a certificate issued earlier could remain legally valid at the time of listing, so a
seller could hold and display a valid certificate that the registry cannot corroborate. The data
bear this out. The share of listings displaying a score with no registry record falls monotonically
as the registry accumulates, from $62.5\%$ of apartment listings in 2021Q3 to $18.4\%$ in 2024Q4,
mirrored by a rise in verified disclosers from $7.4\%$ to roughly one half.

This coverage gap has a known direction. A seller holding a still-valid pre-reform certificate who
withholds it is recorded as genuinely uninformed rather than as a strategic hider, so the measured
hiding rate $\hat\sigma_{0,m}$ understates the true one. Because the equilibrium belief
$e(\sigma_0)$ is decreasing in $\sigma_0$, an understated hiding rate implies an overstated belief
and hence an overstated concealment cost. The coverage gap therefore reinforces, rather than
threatens, the interpretation of $\hat c_{0,m}$ as an upper bound. It is also one reason the
structural sample begins in 2022, since the second half of 2021, where the gap is most severe, is
excluded.

\subsection{Stylized Facts}
\label{sec:facts}

Four features of the market discipline any candidate theory. Each is documented for both segments,
apartments first.

\subsubsection{Non-Disclosing Listings Command a Price Premium, Widened by the 2021 Reform}
\label{ss:fact1}

The starting puzzle is that silence about energy performance commands a price premium relative to the worst certified scores in both property types. Figure~\ref{fig:price_levels} plots the volume-weighted average of market-level hedonic price indices per square metre for each EPC disclosure category from January 2019 to June 2024, where each index holds property characteristics constant within a local market so that comparisons across disclosure categories reflect differences in the disclosure decision rather than differences in property type or location. For apartments (panel~(a)), non-disclosing listings occupy the highest price tier across the \emph{entire} sample window: the ordering Non-disclosed $>$ Efficient (A/B/C) $>$ Energy-intensive (E/F/G) is visible as early as 2019, well before the reform, and the gap widens sharply after July 2021. For houses (panel~(b)), the pre-reform picture differs: non-disclosers begin near the bottom of the price distribution and rise toward the middle post-reform, reflecting the higher baseline audit coverage in the house segment documented in Section~\ref{ss:fact3}. The pre-reform level premium for apartments is consistent with the concentration of non-audited stock in high-price urban markets. Figure~\ref{fig:fact1_premium} isolates the \emph{change} in the gap: it plots difference-in-differences estimates of the hedonic-controlled log price gap between non-disclosing and energy-intensive listings, purging observable composition differences and measured relative to the pre-reform baseline. Because the estimates are normalised to that baseline, the pre-reform coefficients are informative about trends rather than levels. They average $+0.002$ log points for apartments and $+0.008$ for houses and display no discernible drift, indicating that the two series evolved in parallel before July 2021. This is precisely the identifying assumption the design requires; it is not a claim that non-disclosing and energy-intensive listings were priced alike, and Figure~\ref{fig:price_levels} shows that for apartments they were not. After July 2021, when DPE certificates became legally opposable under \emph{loi Climat et R\'esilience}, the estimates shift abruptly upward and stabilize around $+0.044$ log points for apartments and $+0.037$ log points for houses by mid-2024. The reform therefore does not create the level difference so much as widen it sharply, in both the raw series and the composition-adjusted estimates. Classical unraveling predicts the opposite sign: silence should signal bad news and trade \emph{below} verified low-quality listings.

\begin{figure}[tbp]
    \centering
    \caption{Price per m\textsuperscript{2} by EPC disclosure category}
    \label{fig:price_levels}
    \begin{subfigure}{0.49\textwidth}
        \centering
        \caption{Apartments}
        \includegraphics[width=\textwidth]{fig/stylized_facts/apartment/fig_price_levels.pdf}
    \end{subfigure}
    \begin{subfigure}{0.49\textwidth}
        \centering
        \caption{Houses}
        \includegraphics[width=\textwidth]{fig/stylized_facts/house/fig_price_levels.pdf}
    \end{subfigure}
    \begin{minipage}{0.95\textwidth}
        \footnotesize\textit{Notes.} Volume-weighted national average of market-level hedonic price indices per square metre (\euro) by EPC disclosure category, January 2019 to June 2024. Price indices are constructed in \citet{patras_mandatory_2024} using hedonic regressions estimated separately within each \emph{zone d'emploi} (INSEE clusters of municipalities), covering 110 local markets. National series are obtained by aggregating market-level indices by listing volume. Series are smoothed with a centred three-month moving average. The vertical dashed line marks the July 2021 DPE reform. Source: \citet{patras_mandatory_2024}; authors' calculations.
    \end{minipage}
\end{figure}

\begin{figure}[tbp]
    \centering
    \caption{Non-disclosure price premium around the July 2021 DPE reform}
    \label{fig:fact1_premium}
    \begin{subfigure}{0.49\textwidth}
        \centering
        \caption{Apartments}
        \includegraphics[width=\textwidth]{fig/stylized_facts/apartment/fig_fact1_nondiscl_premium_eventstudy.pdf}
    \end{subfigure}
    \begin{subfigure}{0.49\textwidth}
        \centering
        \caption{Houses}
        \includegraphics[width=\textwidth]{fig/stylized_facts/house/fig_fact1_nondiscl_premium_eventstudy.pdf}
    \end{subfigure}
    \begin{minipage}{0.95\textwidth}
        \footnotesize\textit{Notes.} Monthly volume-weighted log gap between the price index of non-disclosing listings and certified energy-intensive (E/F/G) listings, January 2019 to June 2024. The vertical dashed line marks the legal opposability of the DPE (July 2021). The thick line is a centred three-month moving average. Source: SeLoger residential listings matched against the ADEME registry of certified energy audits; authors' calculations.
    \end{minipage}
\end{figure}

\subsubsection{Non-Disclosure Persists Long After Mandatory Disclosure}
\label{ss:fact2}

Classical unraveling further predicts that, once disclosure is mandatory and verifiable, the share of non-disclosing listings should converge rapidly toward zero. The French data reject this prediction in both segments, though with different magnitudes. Figure~\ref{fig:fact2_missing} plots the monthly share of listings without a certified EPC from October 2018 through July 2024. For apartments (panel~(a)), the share averages $35.9\%$ before the reform and $33.8\%$ afterward; over the final twelve months of the sample, it stabilizes at $28.9\%$. For houses (panel~(b)), the corresponding numbers are $24.6\%$ pre-reform, $14.4\%$ post-reform, and a final-twelve-month mean of $9.7\%$. Houses thus display better baseline audit coverage and faster post-reform convergence, but in both segments the share remains structurally bounded away from zero more than three years after disclosure became legally required. Almost three-tenths of apartment listings and one-tenth of house listings still come to market without a certified audit, indicating that a substantial fraction of sellers cannot disclose, because they do not possess a certified audit at the time of listing.

\begin{figure}[htbp]
    \centering
    \caption{Persistence of non-disclosure after the DPE reform}
    \label{fig:fact2_missing}
    \begin{subfigure}{0.49\textwidth}
        \centering
        \caption{Apartments}
        \includegraphics[width=\textwidth]{fig/stylized_facts/apartment/fig_fact2_missing_share_timeseries.pdf}
    \end{subfigure}
    \begin{subfigure}{0.49\textwidth}
        \centering
        \caption{Houses}
        \includegraphics[width=\textwidth]{fig/stylized_facts/house/fig_fact2_missing_share_timeseries.pdf}
    \end{subfigure}
    \begin{minipage}{0.95\textwidth}
        \footnotesize\textit{Notes.} Monthly share of listings without a certified EPC, October 2018 to July 2024. The horizontal dashed line marks each segment's final-twelve-month mean ($0.29$ apartments, $0.10$ houses). The vertical dashed line marks the July 2021 reform date. Volumes are summed across markets within each segment. Source: SeLoger residential listings matched against the ADEME registry of certified energy audits; authors' calculations.
    \end{minipage}
\end{figure}

\subsubsection{Direct Decomposition of the Listing Pool by Information Status}
\label{ss:fact3}

A central methodological contribution of this paper is that, by matching every listing against the ADEME registry, we can directly observe whether a seller is informed ($K=1$, audit exists) or uninformed ($K=0$, no audit on record). Crossing this latent information status with the seller's disclosure choice $d \in \{0,1\}$ yields a four-cell classification of which only three cells are economically distinct: truthful disclosers ($K=1, d=1$), strategically hiding sellers ($K=1, d=0$), and genuinely uninformed sellers ($K=0, d=0$).\footnote{The fourth cell ($K=0$, $d=1$) comprises sellers who display an EPC score for which no matching record appears in the public ADEME registry. These listings are treated as truthful disclosers, consistent with the no-falsification assumption: a score published by the owner is taken at face value, and only omission is treated as a strategic act. Two institutional features explain the size of this group. First, the public registry contains only certificates issued from July 2021 onward, whereas a certificate issued earlier could still be legally valid at the time of listing. Under \emph{d\'ecret} n\textsuperscript{o}~2020-1610 of 17 December 2020, certificates issued between 1 January 2013 and 31 December 2017 remained valid until 31 December 2022, and those issued between 1 January 2018 and 30 June 2021 until 31 December 2024, while certificates issued from 1 July 2021 are valid for ten years. A seller listing during our sample window could therefore display a still-valid certificate that is absent from the registry by construction. Second, the model assumes no falsification. Consistent with the first explanation, the share of listings in this cell falls monotonically as the registry accumulates post-reform certificates: from $62.5\%$ of apartment listings in 2021Q3 to $18.4\%$ in 2024Q4, and from $73.1\%$ to $43.7\%$ for houses.} Figure~\ref{fig:fact3_decomposition} reports the resulting decomposition for both segments. In apartments (panel~(a)), $79.8\%$ of listings are truthful disclosers, $16.7\%$ are genuinely uninformed, and $3.5\%$ are strategically hiding. In houses (panel~(b)), the corresponding shares are $90.4\%$, $8.5\%$, and $1.1\%$. Houses systematically display higher baseline disclosure and a smaller strategic-hiding share, but the qualitative structure is identical: a non-trivial mass of strategic hiders coexists with a much larger mass of genuinely uninformed sellers in both segments.

\begin{figure}[htbp]
    \centering
    \caption{Three-group decomposition of the listing pool}
    \label{fig:fact3_decomposition}
    \begin{subfigure}{0.49\textwidth}
        \centering
        \caption{Apartments}
        \includegraphics[width=\textwidth]{fig/stylized_facts/apartment/fig_fact3_three_group_decomposition.pdf}
    \end{subfigure}
    \begin{subfigure}{0.49\textwidth}
        \centering
        \caption{Houses}
        \includegraphics[width=\textwidth]{fig/stylized_facts/house/fig_fact3_three_group_decomposition.pdf}
    \end{subfigure}
    \begin{minipage}{0.95\textwidth}
        \footnotesize\textit{Notes.} Shares from the estimation sample of each property type: 834{,}816 apartment listings and 677{,}011 house listings, after registry-match cleaning and market-size filtering. Source: SeLoger residential listings matched against the ADEME registry of certified energy audits; authors' calculations.
    \end{minipage}
\end{figure}

\noindent Two features of this decomposition are central to the analysis. First, the non-disclosing pool ($d=0$) is mechanically composite in both segments: roughly five-sixths of apartment non-disclosers and seven-eighths of house non-disclosers are genuinely uninformed, while the remainder are strategically concealing. The vast majority of silent listings are silent by constraint, not by choice. Second, the relative weight of these two subgroups varies substantially across property types and, as we show below, across local markets, with direct implications for the equilibrium informativeness of non-disclosure.

\subsubsection{Substantial Cross-Market Heterogeneity in Non-Disclosure}
\label{ss:fact4}

Figure~\ref{fig:fact4_map} maps the apartment (panel~(a)) and house (panel~(b)) non-disclosure rates at the d\'epartement level. For apartments, the rate ranges from $6.5\%$ in Aube to $40.9\%$ in Corse-du-Sud, with a cross-d\'epartement median of $16.5\%$; high values cluster in southern France, Corsica, and parts of the Parisian region (Seine-Saint-Denis at $32.1\%$). For houses, the range is $5.4\%$ in Essonne to $37.4\%$ in Paris, with a median of $10.0\%$; the Paris d\'epartement is an outlier on the high end, while southern and rural d\'epartements lie closer to the metropolitan baseline. The geographic patterns differ across the two segments, reflecting the different local composition of each property type's listing pool, but the underlying message is identical: within-segment cross-d\'epartement dispersion is large in both cases. This heterogeneity rules out a national-level treatment of the information environment in either segment: any credible structural exercise must allow the depth of audit coverage and the salience of EPC information to vary by local market.

\begin{figure}[htbp]
    \centering
    \caption{Non-disclosure rate by d\'epartement, July 2021 to December 2024}
    \label{fig:fact4_map}
    \begin{subfigure}{0.49\textwidth}
        \centering
        \caption{Apartments}
        \includegraphics[width=\textwidth]{fig/stylized_facts/apartment/fig_fact5_nondiscl_rate_map.pdf}
    \end{subfigure}
    \begin{subfigure}{0.49\textwidth}
        \centering
        \caption{Houses}
        \includegraphics[width=\textwidth]{fig/stylized_facts/house/fig_fact5_nondiscl_rate_map.pdf}
    \end{subfigure}
    \begin{minipage}{0.95\textwidth}
        \footnotesize\textit{Notes.} Shares are computed on listings with a valid d\'epartement code; d\'epartements with fewer than 200 listings within each segment are shown in light grey, and overseas d\'epartements are omitted. Classification: quantiles ($k=5$) on metropolitan France, computed separately within each panel so the breakpoints reflect each segment's distribution. Projection: Lambert-93 (EPSG:2154). Source: SeLoger residential listings matched against the ADEME registry of certified energy audits; d\'epartement boundaries from \texttt{france-geojson} (G.~David, 2024).
    \end{minipage}
\end{figure}

\subsection{What the Facts Demand of the Theory}
\label{ss:facts_demand}

Taken together, the four stylized facts impose tight discipline on any candidate theory of EPC disclosure. The documented non-disclosure premium rules out classical full-unraveling models with skeptical buyers, in which silence must trade at or below the worst certified score. The persistence of non-disclosure after the reform rules out frictionless implementations in which enforcement is complete and disclosure is universal at any positive cost of concealment. The composite nature of the non-disclosing pool requires the theory to recognize that silence aggregates two economically distinct subpopulations: informed sellers who choose to conceal unfavorable information, and uninformed sellers who cannot disclose because no certified audit exists. Their joint mass determines the informational content of non-disclosure and cannot be collapsed into a single type. Finally, the substantial cross-market heterogeneity in non-disclosure rates and price gaps rules out a uniform national treatment of these structural parameters: estimation must be conducted market by market. The model developed in the next section is designed to satisfy these four requirements simultaneously.
